\chapter{Supramoleculaire chemie}\label{ch:b3:supramolecular}

In 1962 isoleerde Charles Pedersen, die naar iets anders zocht, een
cyclische polyether die kaliumpermanganaat in benzeen oploste en het
paars kleurde. De ring had zich rond het kaliumion gewikkeld en zijn
lading achter een koolwaterstofoppervlak verborgen, en sleepte het
permanganaat als tegenion mee. De chemie voorbij het molecuul, de chemie
van assemblages die bijeengehouden worden door krachten die zwakker zijn
dan covalente bindingen, begon met die paarse oplossing.
Dit hoofdstuk meet die krachten, legt uit waarom sommige gastheren hun
\omterm{def:b3:supramolecular:supramolecular}{gasten} zo goed kiezen, toont hoe bindingsconstanten gemeten worden, en
eindigt met moleculen die tot machines geregen en vervlochten zijn.

\begin{recall}
Het deel van jaar~1 behandelde intermoleculaire krachten, de
waterstofbrug, solvatatie, hydrofobe moleculen, complexen,
vormingsconstanten en chelaten; het deel van jaar~2 de chemische
potentiaal, gibbsenergieën en de kleinstekwadratenaanpassing.
\Cref{ch:b3:advanced-nmr} beschreef snelle en trage uitwisseling en
coalescentie,
\cref{ch:b3:rate-theories} de \omterm{thm:b3:rate-theories:eyring}{vergelijking van Eyring}, en
\cref{ch:b3:partition-functions} de translatie-entropie van een gas.
\end{recall}

\section{Niet-covalente interacties}

\begin{definition}[Supramoleculaire chemie]\label{def:b3:supramolecular:supramolecular}
\emph{Supramoleculaire chemie}\index{supramoleculaire chemie} is de chemie
van assemblages van moleculen die bijeengehouden worden door
niet-covalente interacties. In een
\emph{gastheer-gastcomplex}\index{gastheer-gastcomplex} is de
\emph{gastheer}\index{gastheer} het grotere molecuul, met een holte of een
spleet en bindingsplaatsen die erop gericht zijn, en de
\emph{gast}\index{gast} het kleinere molecuul of ion dat hij bindt.
\end{definition}

De interacties zijn die van het deel van jaar~1, ion-dipool,
dipool-dipool, waterstofbruggen en londonkrachten, plus enkele met een
eigen naam.

\begin{definition}[Niet-covalente interacties]\label{def:b3:supramolecular:interactions}
\emph{$\pi$-stapeling}\index{$\pi$-stapeling} is de aantrekking tussen vlakke
aromatische ringen die met hun vlakken tegenover elkaar liggen, meestal
verschoven of gekanteld. Een
\emph{kation-$\pi$-interactie}\index{kation-$\pi$-interactie} is de
aantrekking tussen een kation en het elektronenrijke vlak van een
aromatische ring. Een \emph{halogeenbinding}\index{halogeenbinding}
is de aantrekking tussen de elektronenarme punt van een covalent gebonden
halogeenatoom, tegenover zijn binding, en een lewisbase. Het
\emph{hydrofoob effect}\index{hydrofoob effect} is de neiging van apolaire
oppervlakken om in water samen te komen, vooral gedreven door het
vrijkomen van de geordende watermoleculen eromheen.
\end{definition}

De ion-dipool- en kation-$\pi$-interacties zijn de sterkste, en hangen sterk af
van het oplosmiddel dat om het ion concurreert; daarna komen de
waterstofbruggen, dan de $\pi$-stapeling en de \omterm{def:b3:supramolecular:interactions}{halogeenbindingen}, dan de
londonkrachten, die afzonderlijk zwak zijn maar optellen over grote
contactoppervlakken. In water overheerst vaak het \omterm{def:b3:supramolecular:interactions}{hydrofoob effect}: een
\omterm{def:b3:supramolecular:supramolecular}{gastheer} die ontworpen is om een \omterm{def:b3:supramolecular:supramolecular}{gast} in chloroform via waterstofbruggen
te binden, bindt hem in water misschien helemaal niet, omdat de eigen
waterstofbruggen van het water concurreren.

\section{Moleculaire herkenning}

\begin{definition}[Moleculaire herkenning]\label{def:b3:supramolecular:recognition}
\emph{Moleculaire herkenning}\index{moleculaire herkenning} is de
selectieve binding van een \omterm{def:b3:supramolecular:supramolecular}{gast} door een \omterm{def:b3:supramolecular:supramolecular}{gastheer}. Ze berust op
\emph{complementariteit}\index{complementariteit}, een overeenstemming
van grootte, vorm en bindingsplaatsen tussen \omterm{def:b3:supramolecular:supramolecular}{gastheer} en \omterm{def:b3:supramolecular:supramolecular}{gast}, en wordt
versterkt door \emph{preorganisatie}\index{preorganisatie}: een \omterm{def:b3:supramolecular:supramolecular}{gastheer}
waarvan de bindingsplaatsen al vóór de binding in de geometrie van het
complex gehouden worden.
\end{definition}

De associatieconstante $K_a = [\mathrm{HG}]/([\mathrm H][\mathrm G])$ is een vormingsconstante in de zin van
het deel van jaar~1, en $\Delta G^\circ = -RT\ln K_a$. Een flexibele \omterm{def:b3:supramolecular:supramolecular}{gastheer} moet bij de
binding zijn conformatie bevriezen, wat entropie kost; een starre,
gepreorganiseerde \omterm{def:b3:supramolecular:supramolecular}{gastheer} heeft die prijs al bij zijn synthese betaald.

\begin{definition}[Chelaat- en macrocyclisch effect]\label{def:b3:supramolecular:macrocyclic}
Het \emph{chelaateffect}\index{chelaateffect} is de grotere stabiliteit van
een complex met een meertandig ligand dan met het overeenkomstige aantal
eentandige liganden met dezelfde donoratomen. Het
\emph{macrocyclisch effect}\index{macrocyclisch effect} is de bijkomende
stabilisatie wanneer het meertandige ligand een ring is in plaats van een
open keten.
\end{definition}

\begin{proposition}[Het chelaateffect is vooral entropisch]\label{prop:b3:supramolecular:chelate-entropy}
In $\mathrm{[M(NH_3)_6]^{2+} + 3\,en \rightleftharpoons [M(en)_3]^{2+} + 6\,NH_3}$, waarin en ethaan-1,2-diamine is, zijn aan beide kanten
dezelfde zes metaal-stikstofbindingen aanwezig, en de reactie wordt
gedreven door de toename van het aantal vrije moleculen.
\end{proposition}

\begin{proof}[Beargumenteerd]
De gevormde en verbroken bindingen zijn van hetzelfde soort, dus $\Delta_rH^\circ$ is klein. Het aantal
onafhankelijke moleculen gaat van vier naar zeven: drie translatievrijheidsgraden
van een heel molecuul meer, elk tientallen joule per kelvin
en per mol waard in oplossing (\cref{ch:b3:partition-functions}), dus $\Delta_rS^\circ > 0$ en $K > 1$.
Kinetisch bekeken: zodra één uiteinde van een chelaat gebonden is, wordt
het andere bij een hoge effectieve concentratie dicht bij het metaal
gehouden, en bindt het meteen opnieuw als het loslaat.
\end{proof}

\begin{proposition}[Enthalpie-entropiecompensatie]\label{prop:b3:supramolecular:compensation}
In een reeks verwante \omterm{def:b3:supramolecular:supramolecular}{gastheer-gastcomplexen} gaat een negatievere $\Delta H^\circ$ van binding
meestal samen met een negatievere $\Delta S^\circ$, zodat $\Delta G^\circ$ minder varieert dan elk van beide (een
experimentele waarneming).
\end{proposition}

Een strakker complex bindt sterker, maar beperkt de bewegingen van
\omterm{def:b3:supramolecular:supramolecular}{gastheer} en \omterm{def:b3:supramolecular:supramolecular}{gast} meer; een sterkere waterstofbrug met de \omterm{def:b3:supramolecular:supramolecular}{gast} betekent
vaak een zwakkere met het water dat hij verdrong. De gibbsenergie, die de
selectiviteit bepaalt, is het kleine verschil van twee grote termen.

\begin{definition}[Gastheren]\label{def:b3:supramolecular:hosts}
Een \emph{kroonether}\index{kroonether} is een macrocyclische polyether,
zoals
18-kroon-6, $\ce{(CH2CH2O)6}$. Een \emph{cryptand}\index{cryptand} is een
bicyclische polyether met twee bruggenhoofden van stikstof, die een
kation in drie dimensies omsluit in een \emph{cryptaat}\index{cryptaat}.
Een \emph{cyclodextrine}\index{cyclodextrine} is een cyclisch oligomeer
van glucose-eenheden, in de vorm van een afgeknotte kegel met een
hydrofobe holte en hydroxygroepen op zijn randen. Een
\emph{calixareen}\index{calixareen} is een bekervormige macrocyclus van
fenoleenheden die door methyleenbruggen verbonden zijn.
\end{definition}

\begin{omfigure}
\begin{tikzpicture}[font=\small]
  % 18-crown-6 with K+
  \foreach \k in {0,...,17} {
    \pgfmathsetmacro\ang{90 + 20*\k}
    \pgfmathsetmacro\nxt{110 + 20*\k}
    \draw[black!70] (\ang:1.55) -- (\nxt:1.55);
  }
  \foreach \k in {0,3,...,15} {
    \pgfmathsetmacro\ang{90 + 20*\k}
    \node[fill=white, inner sep=0.5pt, text=omThm] (o\k) at (\ang:1.55) {O};
    \draw[omThm!60, densely dashed] (0,0) -- (\ang:1.32);
  }
  \node[circle, fill=omDef!25, draw=omDef, inner sep=1pt] at (0,0) {$\ce{K+}$};
  \node[font=\scriptsize] at (0,-2.1) {18-kroon-6 met $\ce{K+}$};
  % [2.2.2]cryptand with K+
  \begin{scope}[xshift=5.2cm]
    \node (nt) at (0,1.7) {N};
    \node (nb) at (0,-1.7) {N};
    \foreach \x/\n in {-1.25/a, 0/b, 1.25/c} {
      \node[text=omThm, fill=white, inner sep=0.5pt] (oa\n) at (\x,0.55) {O};
      \node[text=omThm, fill=white, inner sep=0.5pt] (ob\n) at (\x,-0.55) {O};
      \draw[black!70] (nt) to[out=-90+40*\x, in=90] (oa\n) -- (ob\n) to[out=-90, in=90-40*\x] (nb);
    }
    \node[circle, fill=omDef!25, draw=omDef, inner sep=1pt] at (0.62,0) {$\ce{K+}$};
    \node[font=\scriptsize] at (0,-2.4) {[2.2.2]cryptand: cryptaat van $\ce{K+}$};
  \end{scope}
\end{tikzpicture}

{\small Links: 18-kroon-6, een ring van twaalf koolstofatomen
(hoekpunten) en zes zuurstofatomen, met een kaliumion in het centrum,
gebonden door de zes zuurstofatomen. Rechts: de
[2.2.2]\omterm{def:b3:supramolecular:hosts}{cryptand}, drie ketens met elk twee etherzuurstofatomen tussen twee
bruggenhoofden van stikstof, schematisch: het kation wordt in een
driedimensionale kooi gehouden door zes zuurstofatomen en de twee
stikstofatomen.}
\end{omfigure}

\omterm{def:b3:supramolecular:hosts}{Kroonethers} selecteren kationen op grootte: 18-kroon-6 bindt $\ce{K+}$ sterker dan het
kleinere $\ce{Na+}$, dat niet alle zes zuurstofatomen tegelijk kan bereiken, of het grotere $\ce{Cs+}$, dat
boven de ring zit. De \omterm{def:b3:supramolecular:hosts}{cryptanden}, sterker gepreorganiseerd en volledig om
het ion gewikkeld, binden nog sterker en selecteren scherper. Verborgen
in de \omterm{def:b3:supramolecular:supramolecular}{gastheer} sleept het kation zijn anion mee in apolaire
oplosmiddelen, waar het anion, niet gesolvateerd, zeer reactief wordt:
de basis van de fasetransferkatalyse.

\section{Binding meten}

\begin{proposition}[Snelle uitwisseling]\label{prop:b3:supramolecular:fast-exchange}
Als een \omterm{def:b3:supramolecular:supramolecular}{gastheer} veel sneller tussen zijn vrije en zijn gebonden toestand
wisselt dan het verschil van hun resonantiefrequenties, verschijnt zijn
NMR-signaal bij het gemiddelde
$\delta_{\mathrm{obs}} = x_{\mathrm{vrij}}\delta_{\mathrm{vrij}} + x_{\mathrm{geb}}\delta_{\mathrm{geb}}$, waarin $x$ de fracties van de \omterm{def:b3:supramolecular:supramolecular}{gastheer} in elke toestand zijn.
\end{proposition}

\begin{proof}
Elk gastheermolecuul bezoekt tijdens de meting vele keren beide
toestanden; zijn precessiefrequentie, gemiddeld over de tijd, is het
gewogen gemiddelde van de twee, met als gewichten de fracties van de tijd
die het in elke toestand doorbrengt, die in evenwicht gelijk zijn aan de
fracties moleculen in elke toestand (\cref{ch:b3:advanced-nmr}).
\end{proof}

\begin{theorem}[Exacte 1:1-isotherm]\label{thm:b3:supramolecular:isotherm}
Voor $\mathrm{H + G \rightleftharpoons HG}$ met totale concentraties $H_0$ en $G_0$ en associatieconstante $K$ is
\[
  [\mathrm{HG}] = \frac{1}{2}\Bigl[\Bigl(H_0 + G_0 + \frac1K\Bigr) - \sqrt{\Bigl(H_0 + G_0 + \frac1K\Bigr)^2 - 4H_0G_0}\,\Bigr].
\]
\end{theorem}

\begin{proof}
Met $c = [\mathrm{HG}]$ is $K = c/((H_0 - c)(G_0 - c))$, d.w.z.\ $c^2 - (H_0 + G_0 + 1/K)c + H_0G_0 = 0$. Van de twee wortels
ligt alleen de kleinste onder zowel $H_0$ als $G_0$ (hun product is $H_0G_0$ en hun som
is groter dan $H_0 + G_0$): het is die met het minteken.
\end{proof}

\begin{omfigure}
\begin{tikzpicture}
\begin{axis}[name=a, width=0.48\linewidth, height=5.5cm, xmin=0, xmax=6, ymin=0, ymax=1.05,
  axis lines=left, xlabel={equivalenten gast, $G_0/H_0$}, ylabel={gebonden fractie gastheer},
  tick label style={font=\scriptsize}, label style={font=\small},
  ]
\addplot[omMeth, thick] table[x=equiv, y=K4] {figdata/out/bachelor-3-binding-titration.dat};
\addplot[omDef, thick] table[x=equiv, y=K3] {figdata/out/bachelor-3-binding-titration.dat};
\addplot[omThm, thick] table[x=equiv, y=K2] {figdata/out/bachelor-3-binding-titration.dat};
\node[font=\scriptsize, omMeth, anchor=north] at (axis cs:4,0.94) {$K = 10^4$};
\node[font=\scriptsize, omDef, anchor=north] at (axis cs:4,0.68) {$K = 10^3$};
\node[font=\scriptsize, omThm, anchor=north] at (axis cs:4,0.27) {$K = 10^2$};
\end{axis}
\begin{axis}[at={(a.east)}, anchor=west, xshift=1.4cm, width=0.48\linewidth, height=5.5cm,
  xmin=0, xmax=1, ymin=0, ymax=1.05, axis lines=left,
  xlabel={molfractie gastheer}, ylabel={$[\text{complex}]/c_{\text{totaal}}$},
  tick label style={font=\scriptsize}, label style={font=\small}]
\addplot[omDef, thick] table[x=x, y=HG] {figdata/out/bachelor-3-binding-titration-b.dat};
\addplot[omThm, thick] table[x=x, y=HG2] {figdata/out/bachelor-3-binding-titration-b.dat};
\draw[black!35, dashed] (axis cs:0.5,0) -- (axis cs:0.5,0.55) (axis cs:0.3333,0) -- (axis cs:0.3333,0.38);
\node[font=\scriptsize, omDef, anchor=south] at (axis cs:0.5,0.52) {HG};
\node[font=\scriptsize, omThm, anchor=south] at (axis cs:0.25,0.36) {$\mathrm{HG_2}$};
\end{axis}
\end{tikzpicture}

{\small Links: gebonden fractie van een \omterm{def:b3:supramolecular:supramolecular}{gastheer} tijdens een titratie bij $H_0 = \qty{1.0}{mmol/L}$ voor drie
associatieconstanten (in $\unit{L.mol^{-1}}$): hoe groter $K$, hoe scherper de knik bij één equivalent;
als $KH_0 \gg 1$, zegt de kromme alleen dat $K$ groot is. Rechts: grafieken van
Job, de concentratie complex bij een constante totale concentratie van
\omterm{def:b3:supramolecular:supramolecular}{gastheer} plus \omterm{def:b3:supramolecular:supramolecular}{gast}, met een maximum bij een gastheerfractie van 1/2 voor
een complex 1:1 en 1/3 voor een complex 1:2 (modelconstanten).}
\end{omfigure}

\begin{method}[Bindingsconstante uit een NMR-titratie]\label{met:b3:supramolecular:nmr-titration}
\begin{enumerate}
  \item Houd de gastheerconcentratie $H_0$ vast en voeg \omterm{def:b3:supramolecular:supramolecular}{gast} toe, van 0 tot enkele
        equivalenten; registreer bij elk punt een signaal van de \omterm{def:b3:supramolecular:supramolecular}{gastheer}
        (snelle uitwisseling).
  \item Model: $\delta = \delta_{\mathrm{vrij}} + (\delta_{\mathrm{geb}} - \delta_{\mathrm{vrij}})[\mathrm{HG}]/H_0$, met $[\mathrm{HG}]$ uit de exacte isotherm.
  \item Pas $K$ en $\delta_{\mathrm{geb}}$ aan met niet-lineaire kleinste kwadraten, door de som van de
        gekwadrateerde residuen te minimaliseren; haal de standaardfouten
        uit de kromming van die som (de covariantiematrix).
  \item Controleer de residuen op een trend (verkeerde stoichiometrie) en
        of de gebonden fractie ruwweg 20 tot 80\,\% bestrijkt (de gegevens leggen $K$ dan vast).
\end{enumerate}
\end{method}

\begin{definition}[Grafiek van Job]\label{def:b3:supramolecular:job}
Een \emph{grafiek van Job}\index{grafiek van Job} (methode van de continue
variaties) is een grafiek van de concentratie van een complex, of van een
signaal dat ermee evenredig is, tegen de molfractie van één component,
bij een constante totale concentratie van de twee.
\end{definition}

\begin{proposition}[Maximum van een grafiek van Job]\label{prop:b3:supramolecular:job}
Voor één enkel complex $\mathrm{H_nG_m}$ is de \omterm{def:b3:supramolecular:job}{grafiek van Job} maximaal bij de
molfractie \omterm{def:b3:supramolecular:supramolecular}{gastheer}
$x = n/(n + m)$, wat de bindingsconstante ook is.
\end{proposition}

\begin{proof}
Met totale concentratie $T$, $H_t = xT$, $G_t = (1 - x)T$, complex $c$, vrij $[\mathrm H] = H_t - nc$,
$[\mathrm G] = G_t - mc$, en $c = \beta[\mathrm H]^n[\mathrm G]^m$. Afleiden naar $x$ in het maximum, waar
$\dd c/\dd x = 0$: $0 = \beta\bigl(n[\mathrm H]^{n-1}[\mathrm G]^m T - m[\mathrm H]^n[\mathrm G]^{m-1}T\bigr)$, dus $n[\mathrm G] = m[\mathrm H]$. Invullen geeft
$n(G_t - mc) = m(H_t - nc)$, dus $nG_t = mH_t$, d.w.z.\ $n(1 - x) = mx$: $x = n/(n + m)$.
\end{proof}

\begin{method}[Een grafiek van Job opstellen]\label{met:b3:supramolecular:job}
\begin{enumerate}
  \item Maak equimolaire stockoplossingen van \omterm{def:b3:supramolecular:supramolecular}{gastheer} en \omterm{def:b3:supramolecular:supramolecular}{gast}.
  \item Meng ze in verhoudingen $x$ van 0 tot 1 bij een constant totaal volume.
  \item Meet een eigenschap die evenredig is met het complex (een
        absorbantie gecorrigeerd voor de vrije species, of $x\,\Delta\delta$ in NMR).
  \item Lees de stoichiometrie af uit de positie van het maximum.
\end{enumerate}
\end{method}

\begin{definition}[Isotherme titratiecalorimetrie]\label{def:b3:supramolecular:itc}
\emph{Isotherme titratiecalorimetrie}\index{isotherme titratiecalorimetrie}
(ITC) meet de warmte die vrijkomt of opgenomen wordt bij elke injectie van
een gastoplossing in een gastheeroplossing die op constante temperatuur
gehouden wordt.
\end{definition}

\begin{proposition}[Warmte per injectie]\label{prop:b3:supramolecular:itc}
Met verwaarlozing van de verdunning is de warmte van de $i$-de injectie in een cel met volume $V_0$
$q_i = \Delta H^\circ V_0\bigl([\mathrm{HG}]_i - [\mathrm{HG}]_{i-1}\bigr)$.
\end{proposition}

\begin{proof}
De warmte is de reactie-enthalpie maal de hoeveelheid complex die tijdens
de injectie gevormd wordt, en die hoeveelheid is de verandering van $[\mathrm{HG}]$ maal het celvolume.
\end{proof}

\begin{method}[Een ITC-isotherm aflezen]\label{met:b3:supramolecular:itc}
\begin{enumerate}
  \item Zet de warmte per mol geïnjecteerde \omterm{def:b3:supramolecular:supramolecular}{gast} uit tegen de
        molverhouding.
  \item De hoogte van het plateau vóór de verzadiging geeft $\Delta H^\circ$; de molverhouding in
        het buigpunt geeft de stoichiometrie; de steilheid geeft $K$.
  \item Pas alle drie aan met de exacte isotherm; dan is $\Delta G^\circ = -RT\ln K$ en
        $T\Delta S^\circ = \Delta H^\circ - \Delta G^\circ$.
\end{enumerate}
\end{method}

Eén ITC-experiment geeft dus de volledige thermodynamische handtekening
van de binding, enthalpie en entropie, die een NMR- of
uv-vis-titratie alleen via een reeks temperaturen geeft.

\begin{inthelab}[Een NMR-titratie]
Een oplossing van de \omterm{def:b3:supramolecular:supramolecular}{gastheer} in een gedeutereerd oplosmiddel, ongeveer
een millimol per liter, wordt in een NMR-buisje gebracht en haar spectrum
opgenomen. Een geconcentreerde oplossing van de \omterm{def:b3:supramolecular:supramolecular}{gast}, bereid in de
gastheeroplossing zodat de gastheerconcentratie constant blijft, wordt in
kleine porties met een microspuit toegevoegd; na elke toevoeging wordt
het buisje geschud, op de temperatuur van de probe in evenwicht gebracht,
en het spectrum opgenomen. De verschuiving van een signaal van de
\omterm{def:b3:supramolecular:supramolecular}{gastheer} dat beweegt, wordt bij elk punt afgelezen en aangepast.
\end{inthelab}

\section{Gastheren}

\omterm{def:b3:supramolecular:hosts}{Cyclodextrinen}, die door \omterm{def:b3:complex-kinetics:enzyme}{enzymen} uit zetmeel gemaakt worden, hebben zes,
zeven of acht glucose-eenheden ($\alpha$, $\beta$, $\gamma$). Hun hydrofobe holte houdt in water apolaire delen van \omterm{def:b3:supramolecular:supramolecular}{gasten}
vast: ze maken geneesmiddelen oplosbaar, dragen geurstoffen en
verwijderen cholesterol uit celmembranen. Calixarenen en de pompoenvormige
cucurbiturilen, starre macrocycli van glycolurileenheden met openingen
die met carbonylgroepen bekleed zijn, binden kationen en organische
kationen in water zeer sterk.

\begin{omfigure}
\begin{tikzpicture}[font=\small]
  \draw[black!70, fill=omDef!8] (-1.9,1.2) -- (-1.2,-1.2) arc[start angle=180, end angle=360, x radius=1.2, y radius=0.3]
    -- (1.9,1.2);
  \draw[black!70, fill=omDef!18] (0,1.2) ellipse[x radius=1.9, y radius=0.45];
  \draw[black!40, dashed] (-1.2,-1.2) arc[start angle=180, end angle=0, x radius=1.2, y radius=0.3];
  \node[draw=omThm, thick, regular polygon, regular polygon sides=6, minimum size=8mm, inner sep=0pt] at (0,0.2) {};
  \draw[omThm, thick] (0,0.62) -- (0,2.0);
  \node[font=\scriptsize, omThm, anchor=west] at (0.1,2.0) {gast};
  \node[font=\scriptsize, anchor=west] at (2.0,1.2) {brede rand: secundaire OH};
  \node[font=\scriptsize, anchor=west] at (1.3,-1.25) {smalle rand: primaire OH};
  \node[font=\scriptsize, anchor=east] at (-1.75,0) {hydrofobe holte};
\end{tikzpicture}

{\small Een \omterm{def:b3:supramolecular:hosts}{cyclodextrine}, getekend als de afgeknotte kegel die het vormt,
met een aromatische \omterm{def:b3:supramolecular:supramolecular}{gast} in zijn holte. De hydroxygroepen op de twee
randen maken de buitenkant wateroplosbaar; de binnenkant, bekleed met
C--H-groepen en glycosidische zuurstofatomen, is apolair.}
\end{omfigure}

\section{Zelfassemblage en moleculaire machines}

\begin{definition}[Zelfassemblage]\label{def:b3:supramolecular:self-assembly}
\emph{Zelfassemblage}\index{zelfassemblage} is de spontane en omkeerbare
organisatie van componenten tot een bepaalde structuur, gestuurd door de
informatie die in de componenten zelf zit: hun vormen en
bindingsplaatsen.
\end{definition}

Metaalionen met vaste coördinatiehoeken en starre brugliganden
assembleren in één stap tot vierkanten, kooien en capsules, omdat elke
binding kan openen en sluiten tot de stabielste structuur, zonder
spanning of loshangende plaatsen, bereikt is. De basenparing in DNA en de
vouwing van eiwitten zijn hetzelfde idee op grotere schaal.

\begin{definition}[Mechanisch vervlochten moleculen]\label{def:b3:supramolecular:interlocked}
Een \emph{mechanisch vervlochten molecuul}\index{mechanisch vervlochten molecuul}
heeft delen die niet aan elkaar gebonden zijn, maar die niet van elkaar
gescheiden kunnen worden zonder een covalente binding te verbreken. Een
\emph{rotaxaan}\index{rotaxaan} is een ring die op een as geregen is, met
omvangrijke stoppers aan beide uiteinden; een
\emph{catenaan}\index{catenaan} bestaat uit in elkaar grijpende ringen.
\end{definition}

\begin{definition}[Templaatsynthese]\label{def:b3:supramolecular:template}
Een \emph{templaatsynthese}\index{templaatsynthese} gebruikt een tijdelijke
interactie, vaak met een metaalion, om reactieve componenten vast te
houden in de geometrie die nodig is voor een reactie die anders
onwaarschijnlijk zou zijn, zoals het sluiten van een ring rond een andere
ring.
\end{definition}

\begin{omfigure}
\begin{tikzpicture}[font=\small, >=Stealth]
  % step 1: two crescents around Cu+
  \draw[omDef, line width=2pt] (0.3,0) arc[start angle=0, end angle=300, radius=0.75];
  \draw[omThm, line width=2pt] (-0.3,0) arc[start angle=180, end angle=-120, radius=0.75];
  \fill[omProp] (0,0) circle (0.13);
  \node[font=\scriptsize] at (0,-1.35) {$\ce{Cu+}$ houdt twee};
  \node[font=\scriptsize] at (0,-1.7) {open liganden gekruist};
  \draw[->, thick] (1.4,0) -- (2.3,0) node[midway, above, font=\scriptsize] {sluiten};
  % step 2: closed interlocked rings with Cu
  \begin{scope}[xshift=3.6cm]
    \draw[omDef, line width=2pt] (-0.35,0) circle (0.75);
    \draw[omThm, line width=2pt] (0.35,0) circle (0.75);
    \draw[white, line width=5pt] ([shift={(52:0.75)}]-0.35,0) arc[start angle=52, end angle=72, radius=0.75];
    \draw[omDef, line width=2pt] ([shift={(48:0.75)}]-0.35,0) arc[start angle=48, end angle=76, radius=0.75];
    \fill[omProp] (0,0) circle (0.13);
    \node[font=\scriptsize] at (0,-1.35) {catenaat};
  \end{scope}
  \draw[->, thick] (5.0,0) -- (5.9,0) node[midway, above, font=\scriptsize] {$-\ce{Cu+}$};
  \begin{scope}[xshift=7.2cm]
    \draw[omDef, line width=2pt] (-0.35,0) circle (0.75);
    \draw[omThm, line width=2pt] (0.35,0) circle (0.75);
    \draw[white, line width=5pt] ([shift={(52:0.75)}]-0.35,0) arc[start angle=52, end angle=72, radius=0.75];
    \draw[omDef, line width=2pt] ([shift={(48:0.75)}]-0.35,0) arc[start angle=48, end angle=76, radius=0.75];
    \node[font=\scriptsize] at (0,-1.35) {[2]catenaan};
  \end{scope}
  % rotaxane
  \begin{scope}[xshift=10.2cm]
    \draw[black!70, line width=1.6pt] (-1.1,0) -- (1.1,0);
    \fill[black!60] (-1.25,0) circle (0.2) (1.25,0) circle (0.2);
    \draw[omThm, line width=2pt] (0,0) ellipse[x radius=0.18, y radius=0.55];
    \node[font=\scriptsize] at (0,-1.35) {[2]rotaxaan};
  \end{scope}
\end{tikzpicture}

{\small \omterm{def:b3:supramolecular:template}{Templaatsynthese} van een \omterm{def:b3:supramolecular:interlocked}{catenaan} (schematisch): een koper(I)ion
houdt twee liganden op basis van fenantroline loodrecht gekruist in zijn
tetraëdrische coördinatie; het sluiten van elk ervan tot een ring geeft
twee in elkaar grijpende ringen rond het metaal, en het verwijderen van
het koper laat het vrije \omterm{def:b3:supramolecular:interlocked}{catenaan} over. Rechts: een \omterm{def:b3:supramolecular:interlocked}{rotaxaan}, een ring
die door twee stoppers op een as gevangen zit.}
\end{omfigure}

\begin{definition}[Moleculaire machines]\label{def:b3:supramolecular:machine}
Een \emph{moleculaire machine}\index{moleculaire machine} is een
assemblage waarvan de componenten op een gecontroleerde manier ten
opzichte van elkaar bewegen als reactie op een prikkel (licht, een
redoxverandering, een pH-verandering), en die arbeid kan leveren.
\end{definition}

In een moleculaire pendel zit een ring op de as van een \omterm{def:b3:supramolecular:interlocked}{rotaxaan} op een
van twee stations; een prikkel die zijn binding aan het eerste verzwakt,
stuurt hem naar het tweede, en de omgekeerde prikkel stuurt hem terug.
Door licht aangedreven rotatiemotoren op basis van overvolle alkenen
draaien maar in één richting, door fotochemische isomerisaties af te
wisselen met thermische stappen die niet terug kunnen lopen. De snelheid
waarmee de ring tussen de stations pendelt, wordt gemeten met NMR bij
variabele temperatuur, uit de coalescentie van de signalen van de twee
vormen (\cref{ch:b3:advanced-nmr}).

\begin{safety}{\ghs[0.62cm]{GHS07}}
% ledger: ghs4:18C6
18-Kroon-6: schadelijk bij inslikken, irriterend voor de huid, de ogen en
de luchtwegen. \omterm{def:b3:supramolecular:hosts}{Kroonethers} voeren metaalzouten mee in organische
oplosmiddelen en door de huid; ze worden met handschoenen afgewogen en
gehanteerd.
\end{safety}

\begin{history}[De supramoleculaire chemie en haar machines]
De \omterm{def:b3:supramolecular:hosts}{kroonethers} van Charles Pedersen (1967), de \omterm{def:b3:supramolecular:hosts}{cryptanden} van Jean-Marie
Lehn en de gepreorganiseerde gastheren van Donald Cram legden de basis
van het vakgebied; de drie deelden de Nobelprijs voor de Scheikunde van
1987. Jean-Pierre Sauvage maakte in 1983 catenanen met een hoog rendement
via het kopertemplaat, Fraser Stoddart bouwde rotaxaanpendels, en Ben
Feringa in 1999 een door licht aangedreven rotatiemotor; ze deelden de
Nobelprijs voor de Scheikunde van 2016 voor het ontwerp en de synthese
van \omterm{def:b3:supramolecular:machine}{moleculaire machines}.
\end{history}

\section{Oefeningen}

\begin{exercise}[$\star$]\label{exo:b3:supramolecular:1}
Benoem de belangrijkste interactie in: $\ce{K+}$ in 18-kroon-6; twee benzeenringen met de vlakken tegenover elkaar;
$\ce{K+}$ boven een benzeenring; joodpentafluorbenzeen met pyridine; een paar guanine-cytosine.
\end{exercise}

\begin{exercise}[$\star$]\label{exo:b3:supramolecular:2}
Bereken $\Delta G^\circ$ bij \qty{298}{K} voor een \omterm{def:b3:supramolecular:supramolecular}{gastheer-gastcomplex} met $K_a = \qty{1.0e4}{L.mol^{-1}}$.
\end{exercise}

\begin{exercise}[$\star$]\label{exo:b3:supramolecular:3}
In methanol is met 18-kroon-6 $\log K = 6.1$ voor $\ce{K+}$ en 4.4 voor $\ce{Na+}$ (gegevens van de oefening).
Bereken de selectiviteit en het verschil in $\Delta G^\circ$.
\end{exercise}

\begin{exercise}[$\star$]\label{exo:b3:supramolecular:4}
Een \omterm{def:b3:supramolecular:job}{grafiek van Job} heeft een maximum bij een molfractie \omterm{def:b3:supramolecular:supramolecular}{gastheer} van
0.33. Wat is de stoichiometrie?
\end{exercise}

\begin{exercise}[$\star\star$]\label{exo:b3:supramolecular:5}
Een \omterm{def:b3:supramolecular:supramolecular}{gastheer} bij \qty{2.0}{mmol/L} met $K = \qty{1500}{L.mol^{-1}}$, $\delta_{\mathrm{vrij}} = \qty{3.560}{ppm}$ en $\delta_{\mathrm{geb}} = \qty{3.720}{ppm}$ krijgt één
equivalent \omterm{def:b3:supramolecular:supramolecular}{gast}. Bereken de gebonden fractie en de waargenomen
verschuiving.
\end{exercise}

\begin{exercise}[$\star\star$]\label{exo:b3:supramolecular:6}
Voor nikkel(II) is $\log\beta_6 = 8.6$ met ammoniak en $\log\beta_3 = 18.3$ met ethaan-1,2-diamine (gegevens van de
oefening). Bereken de constante en $\Delta G^\circ$ bij \qty{298}{K} van
$\ce{[Ni(NH3)6]^2+} + 3\,\mathrm{en} \rightleftharpoons \ce{[Ni(en)3]^2+} + \ce{6NH3}$, en zeg wat de reactie drijft.
\end{exercise}

\begin{exercise}[$\star\star$]\label{exo:b3:supramolecular:7}
Een ITC-experiment bij \qty{298}{K} geeft $K = \qty{2.0e5}{L.mol^{-1}}$ en $\Delta H^\circ = \qty{-40}{kJ/mol}$. Bereken $\Delta G^\circ$, $T\Delta S^\circ$ en
$\Delta S^\circ$.
\end{exercise}

\begin{exercise}[$\star\star$]\label{exo:b3:supramolecular:8}
Leg de rol uit van koper(I), en van zijn tetraëdrische geometrie, in de
catenaansynthese van Sauvage. Hoe wordt het koper op het einde
verwijderd?
\end{exercise}

\begin{exercise}[$\star\star$]\label{exo:b3:supramolecular:9}
De twee signalen van een rotaxaanpendel, \qty{120}{Hz} uit elkaar, coalesceren bij \qty{300}{K} (gegevens van de
oefening). Bereken de uitwisselingssnelheidsconstante bij de coalescentie
en de \omterm{def:b3:rate-theories:activation}{gibbsenergie van activering} van het pendelen.
\end{exercise}

\begin{exercise}[$\star\star\star$]\label{exo:b3:supramolecular:10}
Een geneesmiddel heeft een intrinsieke oplosbaarheid van \qty{0.10}{mmol/L} en vormt een complex 1:1 met een
\omterm{def:b3:supramolecular:hosts}{cyclodextrine}, $K = \qty{500}{L.mol^{-1}}$ (gegevens van de oefening). Toon aan dat de totale oplosbaarheid in een
oplossing met totale \omterm{def:b3:supramolecular:hosts}{cyclodextrine} $C_t$ gelijk is aan $S_0 + KS_0C_t/(1 + KS_0)$, en bereken haar voor $C_t = \qty{10}{mmol/L}$.
\end{exercise}

\begin{exercise}[$\star\star\star$]\label{exo:b3:supramolecular:11}
Toon met de exacte isotherm aan dat, als $KH_0 \gg 1$, de gebonden fractie gelijk is aan het
aantal equivalenten tot één, en daarna op één blijft. Waarom geeft zo'n
titratie alleen een ondergrens voor $K$?
\end{exercise}

\begin{exercise}[$\star\star\star$]\label{exo:b3:supramolecular:12}
Neem voor de \omterm{def:b3:supramolecular:supramolecular}{gastheer} van \cref{exo:b3:supramolecular:5} aan dat de
associatie \omterm{def:b3:rate-theories:diffusion}{diffusiegecontroleerd} is, $k_{\mathrm{on}} = \qty{1e9}{L.mol^{-1}.s^{-1}}$. Bereken $k_{\mathrm{off}}$ en toon aan dat de uitwisseling snel is op
de NMR-tijdschaal bij \qty{400}{MHz}.
\end{exercise}

\section{Probleem: paarse benzeen}

\begin{problem}[{Weekendprobleem --- paarse benzeen: kalium in een
kroonether, de extractie van permanganaat in een organisch oplosmiddel,
een bindingsconstante uit een NMR-titratie, en waarom een katalytische
hoeveelheid kroonether volstaat}]
\label{pb:b3:supramolecular:1}
% ledger: const:R, ghs:KMnO4, ghs:toluene, rion:K+VI, rion:Na+VI
Gegevens van het probleem. In methanol bij \qty{298}{K} is met 18-kroon-6
$\log K = 6.1$ voor $\ce{K+}$ en 4.4 voor $\ce{Na+}$. Extractie: een waterige fase met \qty{0.10}{mol/L} $\ce{KCl}$ en \qty{1.0}{mmol/L} $\ce{KMnO4}$ wordt geschud
met een gelijk volume tolueen dat de \omterm{def:b3:supramolecular:hosts}{kroonether} L bevat; de enige
geëxtraheerde species is het ionenpaar $\ce{[KL]+}\ce{MnO4-}$, met
$K_{\mathrm{ex}} = [\ce{[KL]+MnO4-}]_{\mathrm{org}}/([\ce{K+}]_{\mathrm{aq}}[\ce{MnO4-}]_{\mathrm{aq}}[\mathrm L]_{\mathrm{org}}) = \qty{1.0e5}{L^2.mol^{-2}}$, en de \omterm{def:b3:supramolecular:hosts}{kroonether} blijft in het tolueen. NMR-titratie van een
\omterm{def:b3:supramolecular:hosts}{kroonether} (\qty{2.0}{mmol/L}) met een natriumzout in gedeutereerd methanol, één signaal $\ce{CH2}$
(ppm) tegen equivalenten $\ce{Na+}$:

\begin{center}\small
\begin{tabular}{l*{11}{c}}
\hline
equiv. & 0 & 0.25 & 0.5 & 0.75 & 1.0 & 1.5 & 2.0 & 3.0 & 4.0 & 6.0 & 8.0\\
$\delta$ & 3.560 & 3.590 & 3.612 & 3.635 & 3.649 & 3.674 & 3.688 & 3.697 & 3.704 & 3.708 & 3.714\\
\hline
\end{tabular}
\end{center}
% table: binding-titration-c

\medskip\noindent\textbf{Deel I --- Het complex.}
\begin{enumerate}
  \item Teken 18-kroon-6 en zijn kaliumcomplex.
  \item Bereken $\Delta G^\circ$ van de binding van $\ce{K+}$ en van $\ce{Na+}$.
  \item Bereken de selectiviteit $\ce{K+}/\ce{Na+}$.
  \item Verklaar de selectiviteit met de zesgecoördineerde stralen $r(\ce{K+}) = \qty{138}{pm}$ en $r(\ce{Na+}) = \qty{102}{pm}$.
  \item Waarom is het complex oplosbaar in tolueen terwijl $\ce{KCl}$ dat niet is?
  \item Waarom bindt de \omterm{def:b3:supramolecular:hosts}{kroonether} zwakker in water dan in methanol?
\end{enumerate}

\medskip\noindent\textbf{Deel II --- Extractie.}
\begin{enumerate}[resume]
  \item Schrijf het extractie-evenwicht op.
  \item Schrijf, met $y$ de geëxtraheerde concentratie, de vergelijking voor $y$ op als
        $[\mathrm L]_0 = \qty{1.0}{mmol/L}$.
  \item Los ze op en geef de geëxtraheerde fractie permanganaat.
  \item Bereken de fractie met $[\mathrm L]_0 = \qty{10}{mmol/L}$.
  \item Bereken haar met $[\mathrm L]_0 = \qty{0.10}{mmol/L}$.
  \item Waarom is het tolueen paars, en waarom is het geëxtraheerde
        permanganaat een sterkere oxidator dan in water?
\end{enumerate}

\medskip\noindent\textbf{Deel III --- Een NMR-titratie.}
\begin{enumerate}[resume]
  \item Waarom beweegt één gemiddeld signaal in plaats van twee signalen?
  \item Schrijf het model voor $\delta$ als functie van de gastconcentratie op.
  \item Pas $K$ en $\delta_{\mathrm{geb}}$ aan met kleinste kwadraten; geef $K$ met zijn standaardfout.
  \item Geef de spreiding van de residuen en becommentarieer.
  \item Over welk gebied van equivalenten ligt de gebonden fractie tussen
        20 en 80\,\%?
  \item Waarom zou dezelfde titratie $K$ voor kalium, $\log K = 6.1$, niet kunnen meten?
\end{enumerate}

\medskip\noindent\textbf{Deel IV --- Fasetransferkatalyse.}
\begin{enumerate}[resume]
  \item Een alkeen in tolueen wordt geoxideerd door het geëxtraheerde
        permanganaat. Waar vindt de reactie plaats?
  \item Wat gebeurt er met de \omterm{def:b3:supramolecular:hosts}{kroonether} na de reactie, en waarom volstaat
        een katalytische hoeveelheid?
  \item Waarom is een overmaat $\ce{KCl}$ in het water nuttig?
  \item Welke goedkopere katalysatoren doen hetzelfde werk in de
        industrie?
  \item Waarom gebruikt men nu tolueen waar Pedersen benzeen gebruikte?
  \item Formuleer het resultaat: de fractie permanganaat die geëxtraheerd
        wordt bij
        $[\mathrm L]_0 = \qty{1.0}{mmol/L}$.
\end{enumerate}
\end{problem}
